% Part: incompleteness
% Chapter: theories-computability
% Section: computably-axiomatizable

\documentclass[../../../include/open-logic-section]{subfiles}

\begin{document}

\olfileid[pt]{inc}{tcp}{cax}

\olsection{Teorias \printtoken{S}{axiomatizable}}

Dizemos que uma teoria $\Th{T}$ é \emph{!!{axiomatizable}} se possui
um conjunto computável de axiomas~$A$. (Dizer que $A$ é um conjunto de
axiomas para $\Th{T}$ significa que
$T = \Setabs{!A}{A \Proves !A}$.) Toda axiomatização ``razoável'' dos
números naturais terá essa propriedade. Em particular, toda teoria com
um conjunto finito de axiomas é !!{axiomatizable}.

\begin{lem}
Suponha que $\Th{T}$ seja !!{axiomatizable}. Então $\Th{T}$ é
!!{computably enumerable}.
\end{lem}

\begin{proof}
Suponha que $A$ seja um conjunto computável de axiomas para $\Th{T}$.
Para determinar se $!A \in T$, basta procurar !!a{derivation} de $!A$
a partir dos axiomas.

Dito de modo um pouco diferente, $!A$ pertence a $\Th{T}$ se, e
somente se, existe uma lista finita de axiomas $!B_1$, \dots, $!B_k$
em $A$ e !!a{derivation} de
$(!B_1 \land \dots \land !B_k) \lif !A$ na lógica de primeira ordem.
Mas já sabemos que todo conjunto com uma definição da forma ``existem
\dots tais que \dots'' é !!{c.e.}, desde que o segundo ``\dots'' seja
computável.
\end{proof}

\end{document}
